\section{The anticanonical cone and a distinguished lc place}
\label{a:cone}

We work with the contradiction sequence introduced after
\eqref{eq:A1}. Thus \(X\to Z\) is birational, \(Z\) is affine,
\(-K_X\) is relatively ample, \(z\in Z\) is closed, and
\(d=\dim X>0\). The integer \(n>1\) is fixed along the sequence and
\(n^{-1}|-nK_X|\) is an lc system. Put \(p=d+1\).
All valuations in this section are trivial on \(k=\C\).

\subsection{Valuations and forms}

We shall use the full initial field of a valuation, as well as its residue
field; these are different objects. We first record the monomial facts
needed to distinguish them. The smooth monomial description and its
finite-ideal refinements are standard in valuation theory; see
\cite[Proposition~3.7 and Lemma~3.6]{JM}. We include the argument in the
form needed for restrictions to subfields.

\begin{lemma}[Monomial charts and restriction]\label{a:abhyankar}
A real valuation of a function field of characteristic zero satisfying
Abhyankar equality is quasi-monomial on a smooth proper model.
The model can dominate a prescribed model and resolve finitely many
additional functions and divisors. Its restriction to a subfunction field
again satisfies Abhyankar equality. In particular, the restriction of a
prime-divisor valuation is either trivial or a positive integral multiple
of a prime-divisor valuation.
\end{lemma}

\begin{proof}
Choose functions with rationally independent values spanning the value
group over \(\Q\), and functions of value zero whose residues form a
transcendence basis of the residue field. Resolve the divisors of the
first functions and the maps to \(\PP^1\) defined by the second, together
with the prescribed data. The latter functions are regular and invertible
at the centre. Consequently the dimension of the centre is at least the
residue transcendence degree. At least as many transverse boundary
components as the rational rank pass through it: the independent values
of the first functions are integral combinations of their orders.
Abhyankar equality forces equality in both bounds. The centre is therefore
the generic point of their transverse intersection, and the normal
parameters have rationally independent values.

In the completed local ring at that generic point, each nonzero function
has a unique least-weight normal monomial with nonzero residue coefficient.
The monomial ideal of strictly higher weight is finitely generated.
The congruence to that monomial times a unit contracts to the local ring
by faithful flatness. This proves the monomial formula.

For a restriction, add the Abhyankar inequalities for the restricted
valuation and the valued field extension. Equality for the full field
forces equality for the restriction. A nonzero subgroup of the value
group \(\Z\) of a prime divisor is \(j\Z\), \(j>0\). The restricted
Abhyankar valuation of rational rank one is therefore \(j\) times a
prime-divisor valuation.
\end{proof}

For a rational top form \(\eta\), its \emph{form discrepancy} is
\(A_\eta(E)=1+\ord_E(\eta)\) at a prime divisor, extended by the
log smooth formula to quasi-monomial valuations. We also write \(a_\eta\).
For a fractional form whose \(r\)-th power is a rational \(r\)-canonical
form, use \(1+r^{-1}\ord_E(\eta^{\otimes r})\).
The top wedge of logarithmic normal differentials and regular tangential
differentials has form discrepancy zero at the monomial valuation;
multiplication by a rational function adds its value.
Pair discrepancy is form discrepancy minus the value of the
log canonical divisor expressed with that form. The trivial valuation
has discrepancy zero. Precise leading-form formulas are proved in
Lemma~\ref{a:leading-form}.

\subsection{The cone and its canonical generator}

Fix a rational volume form \(\omega\) on \(k(X)=k(Z)\) and use
\(K_X=\Div_X(\omega)\). For an indeterminate \(t\), set
\[
 \mathcal A_m=
 \{g t^m:\Div_X(g)-mK_X\geq0\}\cup\{0\},\qquad
 \mathcal A=\bigoplus_{m\geq0}\mathcal A_m,\qquad
 W=\Spec\mathcal A .
\]
Here \(\mathcal A_0=k[Z]\). Denote the degree cocharacter by \(D\), and
the point of the vertex section above \(z\) by \(o\).

\begin{lemma}[Canonical anticanonical cone]\label{a:canonical-cone}
The algebra \(\mathcal A\) is finitely generated and normal, with fraction
field \(k(X)(t)\). The variety \(W\) is klt, and
\[
                   \Omega=t\,\omega\wedge d\log t
\]
has divisor zero on \(W\). In particular it is a canonical generator,
of \(D\)-weight one.
\end{lemma}

\begin{proof}
Relative ampleness gives finite generation of a sufficiently divisible
Veronese and finite generation of each shifted-degree module over it.
Thus it gives finite generation in all degrees. The ring is the
intersection of \(k(X)[t]\) with the valuation rings imposing
\(\ord_P(g)-mP(K_X)\geq0\) at strict prime divisors \(P\) of \(X\).
Here and below, a strict prime of a specified model means a prime divisor
on that model itself; a prime over the model may lie on any higher
birational model. Thus a divisor extracted on a new model is strict on
that model. The displayed intersection is therefore normal. Over the birational generic point of \(Z\)
its field is \(k(X)(t)\).

Consider
\[
 \widehat W=\Spec_X\bigoplus_{m\geq0}\OO_X(-mK_X)t^m .
\]
For a sufficiently divisible \(r\), the analogous relative spectrum in
degrees divisible by \(r\) is the total space of a line bundle. It is
proper over the Veronese cone: evaluation by generating sections embeds
it, relative to \(X\), into a product with affine space over \(Z\), and
\(X\to Z\) is proper. The original relative spectrum is finite over this
line-bundle model. Since \(W\) is separated over its Veronese cone, the
resulting map \(\widehat W\to W\) is proper. It is birational by the
fraction-field calculation. The same prime inequalities give normality
of \(\widehat W\).

Over the smooth locus of \(X\), if \(s\) is a local rational function with
\(\Div(s)=K_X\), the regular fibre coordinate is \(st\), and the displayed
form is a nowhere-vanishing top form. The complement of this locus in
\(\widehat W\) has codimension at least two, as is also seen by the finite
map to the line-bundle model.

For completeness, the local singularities of \(\widehat W\) are klt.
Trivialize \(rK_X\) locally and take the normalized cyclic cover adjoining
\(s\) with \(s^r\) equal to its rational trivializing function.
Because the coefficients of the Weil divisor \(K_X\) are integral,
this cover is quasi-etale. Finite crepant pullback preserves klt
singularities. On the index cover, form the line-bundle algebra
(the normalized reflexive pullback) with coordinate \(y=st\). The deck group acts on \(y\)
with the same character as on \(s\), keeping \(t\) invariant; the form
\(\Omega=(\omega/s)\wedge dy\) is invariant. Its cyclic invariants
recover the original algebra, by the codimension-one regularity inequalities.
The cover of total spaces is again quasi-etale. The finite-cover
discrepancy formula consequently proves that \(\widehat W\) is klt.
We have \(\Div_{\widehat W}(\Omega)=0\), and pushing forward gives
\(\Div_W(\Omega)=0\). Thus the proper birational map is crepant, and
\(W\) is klt as well.
\end{proof}

\begin{lemma}[Gaussian discrepancy]\label{a:cone-gauss}
For a prime divisor \(P\) over \(X\), extend \(P\) by the Gauss rule with
\(t\)-weight \(q-P(K_X)\). If \(q\geq0\), the extension \(\widetilde P\)
is centred on \(W\), and
\[
 \widetilde P(gt^m)=mq+P(\Div_X(g)-mK_X),\qquad
 A_W(\widetilde P)=q+A_X(P).
\]
If \(q\) is rational it is divisorial up to positive scaling.
Every \(D\)-invariant divisorial valuation centred on \(W\) has this form,
allowing a positive multiple of \(P\), or has trivial restriction to
\(k(X)\), in which case its discrepancy is \(q\).
The ideal \(I_n=(\mathcal A_n)\subset\mathcal A\) satisfies
\((W,I_n^{1/n})\) lc.
\end{lemma}

\begin{proof}
The value formula follows from the assigned weight of \(t\).
The divisor \(\Div_X(g)-mK_X\) is effective and rational Cartier,
so the values are nonnegative. A Gauss extension is Abhyankar; wedging
with \(d\log t\) preserves form discrepancy, while multiplication by
\(t\) adds \(q-P(K_X)\). Since
\(A_\omega(P)=A_X(P)+P(K_X)\), this proves the discrepancy formula.
The rational-weight assertion follows from Lemma~\ref{a:abhyankar}.

An invariant valuation evaluates a sum of different torus characters
by their minimum: finite torus translates separate the character
components. It is therefore Gaussian. Properness gives a centre on
\(X\) for its restriction, and Lemma~\ref{a:abhyankar} applies.
The same formulas hold for this restriction and for the trivial one.
The corresponding \(q\) is nonnegative: a sufficiently divisible
relatively free anticanonical section can be chosen with zero divisor
order at the centre, and its value on \(W\) is its degree times \(q\).
The lc assertion reduces to
\[
  \frac1n P\bigl(|-nK_X|\bigr)\leq A_X(P),
\]
which is our choice of \(n\). Invariant divisors suffice: take an
equivariant log resolution of the invariant ideal; the connected torus
fixes its finitely many prime components.

We shall also use general members of these systems. On a log resolution,
the fixed part is the ideal's order and the residual system is free.
A general member with coefficient \(1/n<1\) is transverse to the
fixed support and preserves the lc inequalities. The same choice in a
finite-dimensional generating subspace can therefore be made lc on
both \(X\) and \(W\).
\end{proof}

\subsection{An lc place at the closed vertex}

We need a bounded hypersurface statement only for a fixed-index klt germ.
It is the argument used in \cite[proof of Proposition~5.1]{Chen}.

\begin{lemma}[A hypersurface on a fixed-index germ]\label{a:index-hypersurface}
Given \(p\geq2\) and an integer \(N>0\), there is
\(\tau(p,N)>0\) with the following property. If \((U,G)\) is klt near a
closed point \(x\), \(G\geq0\), \(\dim U=p\), and \(N(K_U+G)\) is
Cartier, then, after shrinking at \(x\), an effective Cartier divisor
\(H_x\) through \(x\) satisfies \((U,G+\tau(p,N)H_x)\) lc.
\end{lemma}

\begin{proof}
Take a plt blowup extracting only \(E\) over \(x\), as in
\cite[Lemma~2.27]{Chen}. For its strict boundary \(G_Y\),
\(-E\) and \(-(K_Y+G_Y+E)\) are relatively ample rational Cartier
divisors. Decreasing the coefficient of \(E\) slightly gives a klt
relative log Fano pair, so \(Y\) is of Fano type over \(U\).
The coefficients of \(G_Y\) belong to the finite \(1/N\)-grid below one.
The relative monotone lc complement theorem
\cite[Theorem~1.8]{BirkarComplements} gives a bounded-index lc complement
of \(G_Y+E\) near \(x\). Its coefficient at \(E\) remains one.

Push the complemented boundary down. The bounded-index principal
identity pushes down as well and makes the pullback crepant, so the
result is an lc boundary \(G^+\geq G\) with \(E\) an lc place.
For a bounded common multiple \(L\) of its complement index and \(N\),
the divisor \(H_x=L(G^+-G)\) is effective Cartier.
It passes through \(x\): its order at \(E\) is
\(L A_{(U,G)}(E)>0\). Taking \(\tau=1/L\), and then a common smaller
positive choice over the finitely many bounded indices, proves the claim.
\end{proof}

In the next proposition, \(N\) is a cutoff tending to infinity along
the contradiction sequence; it is independent of the final complement
index in Theorem~\ref{thm:main}.

\begin{proposition}[A centred lc place detecting the bad sequence]
\label{a:initial-lc-place}
After passing to a subsequence, there are divisorial valuations \(u^0\)
centred at \(o\), and integers \(N\to\infty\), such that
\begin{equation}
 u^0(I_n)=nA_W(u^0),\qquad
 u^0(\m_z)\geq N A_W(u^0).
 \label{eq:A2}
\end{equation}
Here and below \(\m_z\) denotes also its extension to \(\mathcal A\).
\end{proposition}

\begin{proof}
Choose \(N\) divisible by \(n\), tending to infinity slowly enough that
the bad divisors in the contradiction sequence satisfy
\[
 \frac{P(\m_z)}{A_X(P)}>\frac{N}{\tau(p,N)}.
\]
This is a diagonal choice: each fixed \(N\) is allowed before passing
sufficiently far along the bad sequence. For each individual algebra,
choose \(M>N\) sufficiently divisible that the degree-\(M\) elements
cut out the irrelevant locus set-theoretically. Set
\[
                  I=(\mathcal A_N,\m_z,\mathcal A_M).
\]
Its radical is
\(\m_o=\m_z\mathcal A+\mathcal A_{>0}\).
Since \(I_n^{N/n}\subset I\), Lemma~\ref{a:cone-gauss} implies
\((W,I^{1/N})\) lc.

Suppose it were klt near \(o\). A general element \(g\) of a finite
generating space of \(I\) would give
\(G=N^{-1}\Div(g)\) klt there. Indeed the moving part has coefficient
\(1/N<1\) on a log resolution. The canonical generator of \(W\) shows
that \(N(K_W+G)\) is Cartier. Lemma~\ref{a:index-hypersurface} supplies
an effective Cartier hypersurface \(H_o\) through \(o\) for which
\((W,G+\tau(p,N)H_o)\) is lc.

Choose rational \(q\) with
\[
       \frac{A_X(P)}{\tau(p,N)}<q\leq\frac{P(\m_z)}N .
\]
The Gauss extension \(\widetilde P\) has positive value on every
positive-degree element and on \(\m_z\), hence is centred at \(o\).
It has value at least \(Nq\) on \(I\), whereas its bare discrepancy is
\(q+A_X(P)\). Consequently
\[
                  A_{(W,G)}(\widetilde P)\leq A_X(P).
\]
Moreover \(\widetilde P(\m_o)\geq q\): each positive-degree homogeneous
term has value at least its degree times \(q\), and each degree-zero
term vanishing at \(z\) has value at least \(Nq\).
A local equation of \(H_o\) therefore has value at least \(q\).
This contradicts lc, since
\[
 A_{(W,G)}(\widetilde P)\leq A_X(P)
       <\tau(p,N)q\leq\tau(p,N)\widetilde P(H_o).
\]

Thus \(I^{1/N}\) is lc but not klt at \(o\). An lc divisor on a log
resolution is centred there, because \(I\) is supported at \(o\) and
\(W\) is klt elsewhere. For this divisor,
\[
 A_W(u^0)=\frac{u^0(I)}N
       \leq\frac{u^0(I_n)}n\leq A_W(u^0).
\]
Both inequalities are equalities. Since \(\m_z\subset I\), the second
inequality in \eqref{eq:A2} follows too.
\end{proof}
